\subsection{Bibliografia Fundamental}

O PCmin integra a classe de problemas de precificação de múltiplos itens com restrições de demanda. Sua modelagem fundamenta-se em pressupostos sobre o comportamento do consumidor estudados na literatura de otimização combinatória. A formalização geral dos problemas de precificação foi estabelecida por Rusmevichientong et al.~\cite{rusmevichientong2006}, que definiram o cenário onde um vendedor busca maximizar sua receita conhecendo a disposição máxima a pagar de cada cliente. Dentro dessa classe de problemas, o \emph{Problema da Compra Máxima} (PCmax), proposto originalmente por Aggarwal et al.~\cite{aggarwal2004}, adota uma premissa de comportamento inversa ao do PCmin: assume-se que o comprador, diante das opções viáveis, seleciona sempre o item de maior valor absoluto, isto é, se um comprador $b$ adquire o item $i$, a regra de decisão exige que $\pi_i \le v_{ib}$ e que $v_{ib} \ge v_{jb}$ para todo item $j \in I$ tal que $\pi_j \le v_{jb}$. Variando essa dinâmica de escolha, o \emph{Problema da Compra por Preferência} (PCrank) modela cenários onde o consumidor segue uma lista estrita de prioridades pessoais. Embora os critérios de seleção variem entre esses problemas, as estruturas algorítmicas desenvolvidas para eles apoiam indiretamente o avanço do PCmin; a formulação de decisão indireta (DIB), posteriormente desenvolvida por Catabriga~\cite{catabriga2025} para a regra de escolha pelo menor preço, foi originalmente concebida e validada por Aggarwal et al.~\cite{aggarwal2004} no contexto do PCmax.

% Outro problema próximo ao PCmin é o de \emph{Precificação Livre de Inveja} (do inglês, \emph{Envy-Free Pricing}), por Guruswami et al.~\cite{venkatesan2005}. Nessa variante, os preços determinados para os produtos devem garantir um equilíbrio de mercado no qual nenhum comprador prefira a alocação de outro indivíduo à sua própria alocação. O programa linear inteiro misto clássico para o cenário livre de inveja, inicialmente formalizado por Heilporn et al.~\cite{heilporn2010} e expandido por Fernandes et al.~\cite{fernandes2016}, serviu como base direta para Catabriga~\cite{catabriga2025} criar a formulação de Decisão Direta (DD) utilizada no PCmin.

Outro problema próximo ao PCmin é o de \emph{Precificação Livre de Inveja} (do inglês, \emph{Envy-Free Pricing}), introduzido por Guruswami et al.~\cite{venkatesan2005}. Nessa variante, os preços determinados para os produtos devem garantir um equilíbrio de mercado baseado no conceito de utilidade, definida por $u_b = v_{ib} - \pi_i$ se o comprador $b$ recebe o item $i$, e $u_b = 0$ caso contrário, de modo que nenhum comprador prefira a alocação de outro indivíduo à sua própria alocação. Para que a atribuição de um item $i$ ao comprador $b$ seja livre de inveja, deve-se garantir que $v_{ib} - \pi_i \ge 0$ e que $v_{ib} - \pi_i \ge v_{jb} - \pi_j$ para todo produto~$j \in I$, enquanto para um comprador não alocado exige-se que $0 \ge v_{jb} - \pi_j$ para todo $j \in I$. O programa linear inteiro misto clássico para o cenário livre de inveja, inicialmente formalizado por Heilporn et al.~\cite{heilporn2010} através de conexões com a precificação de redes e expandido por Fernandes et al.~\cite{fernandes2016}, serviu como base direta para Catabriga~\cite{catabriga2025} criar a formulação de Decisão Direta (DD) utilizada no PCmin.

Especificamente sobre o avanço teórico do PCmin, Briest e Krysta~\cite{briest2007} demonstraram que o PCmin é NP-difícil. Complementando a análise teórica, Chalermsook et al.~\cite{chalermsook2012} estabeleceram limites de inaproximabilidade para o problema. Ambos os trabalhos evidenciam a elevada dificuldade computacional do problema e justificam o desafio de abordá-lo sob a perspectiva algorítmica. Além disso, recentemente Catabriga~\cite{catabriga2025} realizou o primeiro estudo focado na otimização exata do PCmin, propondo formulações matemáticas inteiras e mistas e investigando técnicas para o fortalecimento de seus limites lineares. Todavia, os experimentos dessa pesquisa pioneira evidenciaram barreiras de escalabilidade para instâncias de grandes dimensões e apontaram uma lacuna completa na literatura no que tange ao desenvolvimento de abordagens heurísticas focadas em cenários de grande porte. Dessa forma, as investigações e desenvolvimentos prévios sobre o PCmax, o PCrank e a precificação livre de inveja fornecem um arsenal metodológico que fundamenta o fortalecimento de modelos e a criação de heurísticas escaláveis propostas neste projeto.

% Uma vez que o Problema da Compra Mínima (PCmin) foi historicamente pouco explorado na literatura sob a ótica da otimização matemática, a fundamentação teórica deste projeto apoia-se fortemente em problemas de precificação correlatos. O PCmax, PCrank e a Precificação Livre de Inveja (Envy-free Pricing) compartilham propriedades estruturais com o PCmin, servindo como alicerce para o avanço das abordagens exatas e aproximadas. 

% Na exploração de métodos exatos, Aggarwal et al. (2004) desenvolveram uma formulação pioneira de PLI para o PCmax no cenário com oferta ilimitada e sem exigência de escada de preços. A estrutura de variáveis e restrições dessa modelagem revelou-se importante na literatura, servindo como base para a formulação de Decisão Indireta Binária (DIB) recentemente adaptada por Catabriga~\cite{catabriga2025} para o PCmin.

% No contexto da Precificação Livre de Inveja, Heilporn et al. (2010) apresentaram uma modelagem de PLIM que associa o problema matemático à Precificação de Rede. Posteriormente, Fernandes et al. (2016) elaboraram novas formulações de PLIM para esse problema, comparando seu desempenho com as abordagens anteriores. Essa lógica de decisão direta, na qual a alocação de itens é definida explicitamente pelas variáveis do modelo, inspirou as modelagens de PLIM (como a de Decisão Direta - DD) desenvolvidas mais tarde também por Catabriga~\cite{catabriga2025} para o PCmin.

% Em relação ao PCrank, Calvete et al. (2019) introduziram formulações não-lineares, apresentaram suas respectivas linearizações e propuseram algoritmos de \emph{branch-and-cut} para a resolução exata. Essa linha de pesquisa foi ampliada por Domínguez et al. (2019), que modelaram o caso de oferta limitada. Adicionalmente, Domínguez, Labbé e Marín formalizaram formulações de PLIM para variantes complexas, como a Compra por Preferência com Empates (2021) e a Compra por Preferência Livre de Inveja (2022). 

% Dado o objetivo deste projeto de desenvolver algoritmos para cenários de grande porte onde métodos exatos geralmente apresentam desempenho computacional comprometido, os avanços em heurísticas para a Precificação Livre de Inveja tornam-se o principal referencial da literatura. Shioda, Tunçel e Myklebust (2011) expandiram heurísticas clássicas para a precificação de múltiplos produtos, combinando-as com cortes válidos para acelerar formulações de PLIM. Aprofundando essa estratégia, Myklebust, Sharpe e Tunçel (2016) propuseram melhorias e implementações eficientes para essas heurísticas, juntamente com novas modelagens exatas. Do ponto de vista teórico, cabe destacar que Aggarwal et al. (2004) desenvolveram um algoritmo de aproximação de razão $O(\log m)$ que pode ser aplicado em conjunto aos problemas PCmax, PCmin e PCrank.

% A validação de melhorias na relaxação linear de modelos PLIM e o teste de novas heurísticas exigem conjuntos de instâncias que simulem o mercado real. Fernandes et al. (2016) introduziram geradores robustos para classes de instâncias esparsas de demanda unitária com bases comportamentais distintas: Characteristics (compradores valoram combinações específicas de características nos itens), Neighborhood (baseado na proximidade em um plano bidimensional) e Popularity (itens populares concentram maior demanda). Shioda, Tunçel e Myklebust (2011) idealizaram também o modelo de instâncias completas (Random Complete), nas quais todos os compradores avaliam todos os itens, que foi posteriormente adotado por Fernandes et al. (2016). Tais instâncias configuram o principal ambiente experimental padrão-ouro adotado na avaliação do desempenho de soluções algorítmicas computacionais para problemas de precificação.

% -----------------------------------------------------

% O $PCmin$ integra uma classe de problemas de precificação de múltiplos itens com restrições de demanda, cujas fundações teóricas e matemáticas apoiam-se em modelos de comportamento do consumidor amplamente estudados na literatura de otimização combinatória. A formalização geral desse ecossistema de precificação dentro de uma abordagem não-paramétrica foi estabelecida por Rusmevichientong et al., que estruturaram o cenário onde um vendedor busca maximizar sua receita conhecendo o conjunto de itens $I = [n]$, o conjunto de compradores $B = [m]$ e a matriz de valorações $v$, onde $v_{ib}$ representa o valor máximo que o comprador $b \in B$ deseja pagar pelo item $i \in I$. O objetivo do vendedor consiste em determinar um vetor de preços $\pi = (\pi_1, \pi_2, \dots, \pi_n)$ e uma matriz binária de alocações $a_{ib}$ que maximizem a receita total $\sum_{b \in B} \sum_{i \in I} \pi_i a_{ib}$. Dentro dessa mesma árvore conceitual, o comportamento do consumidor é modelado por diferentes regras de escolha que geram vertentes distintas. O Problema da Compra Máxima ($PCmax$), proposto originalmente por Aggarwal et al., adota uma premissa de comportamento inversa ao do $PCmin$: assume-se que o comprador, diante das opções viáveis ($\pi_j \le v_{jb}$), seleciona sempre o item de maior valor absoluto. Formalmente, se um comprador $b$ adquire o item $i$, a regra de decisão exige que $\pi_i \le v_{ib}$ e que $v_{ib} \ge v_{jb}$ para todo item $j \in I$ tal que $\pi_j \le v_{jb}$. Variando essa dinâmica de escolha, o Problema da Compra por Preferência ($PCrank$) modela cenários onde o consumidor segue uma lista ou ordem estrita de prioridades pessoais, denotada por $\succ_b$, onde $i \succ_b j$ indica que $b$ prefere o item $i$ ao item $j$. Dessa forma, na alocação do $PCrank$, a compra do item $i$ por um indivíduo $b$ implica que $\pi_i \le v_{ib}$ e que $i \succ_b j$ para todo item viável $j$ tal que $\pi_j \le v_{jb}$.

% Diferente do $PCmin$, que historicamente carecia de abordagens exatas baseadas em programações lineares, o $PCrank$ gerou uma linha de pesquisa robusta focada no desenvolvimento de formulações matemáticas de Programação Linear Inteira Mista (PLIM) e algoritmos de planos de corte. Calvete et al. foram pioneiros ao introduzir formulações para o problema de preferência e demonstrar estratégias de linearização eficazes, propondo os primeiros algoritmos exatos de branch-and-cut fundamentados em pré-processamentos analíticos para a redução do tamanho das instâncias. Esse cenário foi expandido por Domínguez, que propôs uma formulação de PLIM para a variante com capacidades (oferta limitada) e desenvolveu desigualdades válidas adicionadas via branch-and-cut para acelerar a convergência do modelo. Subsequentemente, Domínguez, Labbé e Marín estenderam o realismo dessas estruturas ao formular modelos de PLIM para lidar com o $PCrank$ com Empates na ordem de prioridades e com o $PCrank$ Livre de Inveja, que integra as preferências individuais a critérios de estabilidade e equilíbrio de mercado. Embora os critérios de seleção e a complexidade variem entre essas vertentes, as estruturas algorítmicas desenvolvidas para eles apoiam indiretamente o avanço do $PCmin$; a formulação de decisão indireta (DIB), por exemplo, foi originalmente concebida e validada no contexto do $PCmax$ por Aggarwal et al. antes de ser adaptada para a regra de escolha pelo menor preço.

% Outro pilar teórico essencial para a modelagem matemática do $PCmin$ é o conceito de Precificação Livre de Inveja (do inglês, Envy-Free Pricing), introduzido no campo da otimização algorítmica por Guruswami et al.. Nessa variante, os preços determinados para os produtos devem garantir um equilíbrio de mercado baseado no conceito de utilidade, definida por $u_b = v_{ib} - \pi_i$ se o comprador $b$ recebe o item $i$, e $u_b = 0$ caso contrário. Para que a alocação do item $i$ ao comprador $b$ seja livre de inveja, deve-se garantir formalmente que $v_{ib} - \pi_i \ge 0$ e que $v_{ib} - \pi_i \ge v_{jb} - \pi_j$ para todo produto $j \in I$. Caso o comprador não receba nenhum item, a ausência de inveja exige que $0 \ge v_{jb} - \pi_j$ para todo $j \in I$, assegurando que nenhum indivíduo prefira o feixe alocado a outro em detrimento do seu próprio. O modelo linear inteiro misto clássico para o cenário livre de inveja, inicialmente formalizado por Heilporn et al. através de conexões com o problema de precificação de redes, e posteriormente expandido por Fernandes et al. e Myklebust et al., serve como base direta para a criação da formulação de Decisão Direta (DD) utilizada no $PCmin$.

% Especificamente sobre o avanço do $PCmin$, Briest e Krysta e Chalermsook et al. solidificaram a investigação teórica ao mapear os limites de complexidade e aproximabilidade do problema, demonstrando que ele é NP-difícil e altamente desafiador sob a perspectiva algorítmica em decorrência de fortes limitantes de inaproximabilidade polinomial. Além disso, recentemente Catabriga realizou o primeiro estudo focado na otimização exata do $PCmin$, propondo formulações matemáticas inteiras e mistas ($DIB$, $DD$, $DIV$ e $DDI$) e investigando técnicas baseadas em desigualdades válidas e na redução do espectro de preços para o fortalecimento de seus limites lineares. Todavia, os experimentos dessa pesquisa pioneira evidenciaram sérias barreiras de escalabilidade e memória para instâncias densas ou de grandes dimensões e apontaram uma lacuna completa na literatura no que tange ao desenvolvimento de abordagens heurísticas focadas em cenários de grande porte. Dessa forma, as investigações e desenvolvimentos prévios sobre o $PCmax$, os modelos de PLI para o $PCrank$ e a precificação livre de inveja fornecem o arsenal metodológico que fundamenta o fortalecimento de modelos e a criação de heurísticas escaláveis propostas neste projeto.